\section*{Document control}
\addcontentsline{toc}{section}{Document control}
\begin{longtable}{@{}L{0.22\textwidth}p{0.70\textwidth}@{}}
\toprule
\textbf{Field} & \textbf{Definition} \\
\midrule
Purpose & Define a dependency-ordered path for implementing the core protocols, declarative qualification model, environment-aware resolver, and execution runtime, while keeping concrete implementations in independent packages. \\
Authority & The architecture paper \emph{ImageProcessing: Architecture and Semantic Contracts}, baseline 0.3, is the normative design reference. If this plan conflicts with it, the architecture paper governs until both are revised. \\
Phase completion & Every core phase has required artifacts, tests, evidence, and an explicit exit gate. A partially completed phase remains open. \\
Schedule & No calendar estimates are assigned. Phase size depends on implementation findings, platform access, available hardware, and test evidence. \\
Scope exclusion & No node editor, visual graph frontend, or public graph-programming framework is planned. \\
Implementation principle & Prefer the smallest contract that preserves semantics. The core defines contracts and runtime behavior; applications select and register external provider and operation-implementation packages. \\
\bottomrule
\end{longtable}

\section{Executive roadmap}

\subsection{The governing dependency direction}

The domain API is the architectural priority. \IPP{} defines images, layers, operations, operation specifications, the generic operation-implementation contract, implementation specifications, and declarative qualification-policy types. It consumes generic execution capabilities and environment observations from \GEP{}. The latter is deliberately domain-neutral and MUST NOT import or define image-processing concepts.

The build dependency remains one-way: \IPP{} may import \GEP{}, but \GEP{} must never import \IPP{}. The runtime integrates the two. Concrete execution providers implement \GEP{}; concrete image-operation classes implement \IPP{}. Both kinds of package are selected by the consuming application.

\begin{center}
\begin{tabularx}{0.94\textwidth}{@{}>{\centering\arraybackslash}X@{}}
\fbox{\parbox{0.87\textwidth}{\centering\textbf{Image domain contracts}\quad \IPP{}\\[2pt] Operation semantics, operation specifications, implementation contract/specifications, composition and policy}}\\[6pt]
$\Downarrow$\\[6pt]
\fbox{\parbox{0.87\textwidth}{\centering\textbf{Generic execution contracts}\quad \GEP{}\\[2pt] Providers, devices, resources, static capabilities, availability and state observations}}\\[6pt]
\fbox{\parbox{0.87\textwidth}{\centering\textbf{Runtime and resolver}\quad \IPR{}\\[2pt] Registration, environment snapshots, qualification, routing, planning, execution and result publication}}\\[6pt]
\fbox{\parbox{0.87\textwidth}{\centering\textbf{Application-selected packages}\\[2pt] Provider packs and operation-implementation packs, registered without modifying the core}}
\end{tabularx}
\end{center}

This diagram separates three different things that must not be collapsed:
\begin{itemize}
  \item an \textbf{operation}, such as \texttt{GaussianBlur}, defines what the result means;
  \item an \textbf{operation implementation}, such as \texttt{CUDAGaussianBlur}, implements that operation and declares its static support;
  \item an \textbf{execution provider}, such as CUDA or Metal, supplies generic execution resources and lifecycle.
\end{itemize}

One operation may have multiple implementation classes. Many implementations may use one provider. A managed algorithm library may act as an opaque implementation/runtime adapter rather than as a general provider.

\subsection{Core phase map}

\begin{longtable}{@{}L{0.08\textwidth}L{0.27\textwidth}p{0.54\textwidth}@{}}
\toprule
\textbf{Phase} & \textbf{Name} & \textbf{Completion outcome} \\
\midrule
0 & Architecture decisions and package boundaries & Operation/implementation/environment/policy ownership, numerical dimensions, conflict semantics, environment observability, and package dependencies are decided or explicitly deferred. \\
1 & Image operation and implementation contracts & One typed operation model, one generic implementation protocol, declarative operation/implementation specifications, stable identities and semantic versions. \\
2 & GeneralExecutionProtocol & Generic providers/devices/resources, static capabilities, availability/state observations, telemetry freshness, and asynchronous lifecycle contracts. \\
3 & ImageProcessingProtocol core & Image/layer/function/operation types, composition, normalized fluent policy values, validation, and internal type-erasure boundary. \\
4 & Environment discovery and registries & Provider and implementation discovery/registration, current environment snapshots, and diagnostic provenance. \\
5 & Common qualification and resolution & Shared spec-driven eligibility and ranking using hard requirements, preferences, hardware state, and complete cost estimates. \\
6 & Execution coordination and pipeline planning & Resource reservation/revalidation, residency-aware planning, execution, cancellation, and atomic result publication. \\
7 & Independent implementation packs and reference application & External provider and operation-implementation packages apply the interfaces; a consuming app registers multiple implementations of the same operation without core changes. \\
8 & Measured optimization & End-to-end routing, cost calibration, memory reuse, and safe pipeline optimizations backed by measurements and conformance tests. \\
9 & Production readiness & Versioned distribution, tested support matrix, diagnostics, documentation, consumer integration, and release governance. \\
\bottomrule
\end{longtable}

\subsection{Implementation packages are workstreams, not core phases}

The core roadmap does not schedule CPU, Metal, CUDA, OpenCL, Direct3D 12, Vulkan, OpenCV, or Core Image as sequential core-library phases.

There are two independent implementation-package families:
\begin{enumerate}
  \item \textbf{Execution-provider packages} implement \GEP{} for a technology and expose devices, capabilities, resources, submission, completion, and state observations.
  \item \textbf{Operation-implementation packages} implement one or more \IPP{} operation contracts, publish an \texttt{ImplementationSpecification} for each class, and use an execution provider or managed algorithm runtime.
\end{enumerate}

These packages can be developed independently after their required contracts stabilize. Each is selected by a consuming application according to its target platforms and workloads. A provider package does not automatically provide Gaussian blur, resize, or any other image operation; an operation package does not need to bundle a new provider if it can use an existing one.

\subsection{What is deliberately not a phase}

The roadmap does not include a node editor, public graph editor, visual programming surface, or a second processing engine. Composition is a typed API feature. The runtime may create private execution plans, but those are internal implementation details. Nor does the roadmap include a new neural-network training framework: inference adapters may be added as independent operation implementations when required.

\section{Global rules for phase execution}

\subsection{Definition of done}

A phase is not complete until all applicable conditions hold:
\begin{enumerate}
  \item Its public contracts and decisions are documented.
  \item Its code compiles on every platform claimed by that phase.
  \item Unit and contract tests cover success, invalid input, boundary behavior, and failure paths.
  \item Tests prove that earlier phases' invariants have not regressed.
  \item New provider/operation claims have concrete test evidence on the relevant provider and hardware, not only a mock.
  \item Build and test instructions work from the command line without requiring an IDE-specific project workflow.
  \item Diagnostics identify the operation, implementation, provider/device, and cause when qualification or execution fails.
  \item The README, phase status, and any relevant architecture decision records are updated.
\end{enumerate}

\subsection{Test-first implementation policy}

For each operation, write the semantic reference cases and expected output properties before implementing accelerated kernels. For each provider, write protocol conformance tests before registering it as supported. Tests should be layered:
\begin{itemize}
  \item compile-time/API tests for the public contracts;
  \item pure unit tests for descriptors, parameter validation, identity, and composition;
  \item fake-provider tests for lifecycle, error, and qualification behavior;
  \item CPU reference tests for numerical semantics;
  \item provider-specific tests on actual devices;
  \item end-to-end tests measuring transfers, allocation, synchronization, and complete operation sequences.
\end{itemize}

A fake provider proves that the contract can be exercised; it does not prove GPU correctness. A CPU implementation proves a fallback path; it does not prove a GPU provider. These evidence classes must not be conflated.

\subsection{Compatibility and change control}

An observable change to operation meaning, output interpretation, mutation behavior, or a provider contract requires an explicit version/decision update and regression tests. An implementation optimization that preserves the contract need not change the operation's semantic version. Provider-specific limitations must be declared in implementation specifications and reflected in resolver diagnostics rather than scattered through platform checks in application code.
